% ============================================================
% Chapter 9 — Level 04: Injection Testing
% ============================================================
\chapter{Level 04 — Injection Testing}
\label{ch:level04}

\noteblock{This module is a \textbf{placeholder}. Lab: \texttt{https://sqli.lab.uzyntra.com} (pending deployment).}

\section{Learning Objectives}
\begin{itemize}
  \item Understand SQL Injection concepts and why it occurs
  \item Identify SQL injection points in a web application
  \item Perform manual SQL injection testing with Burp Suite
  \item Understand input validation and secure query patterns
\end{itemize}

\section{Vulnerability Explanation}

SQL Injection occurs when user-supplied input is incorporated into a database query without proper sanitization. The attacker's input changes the logic of the query.

\subsection{Why It Happens}

\begin{lstlisting}[style=code, language=SQL, caption={Vulnerable query construction}]
-- VULNERABLE: string concatenation
query = "SELECT * FROM users WHERE username = '" + username + "'";

-- Attacker input: admin' --
-- Resulting query:
SELECT * FROM users WHERE username = 'admin' --'
-- The '--' comments out the password check!
\end{lstlisting}

\begin{lstlisting}[style=code, language=SQL, caption={Secure parameterized query}]
-- SECURE: parameterized query
query = "SELECT * FROM users WHERE username = ?";
statement.setParameter(1, username);
-- User input is treated as data, never as SQL code
\end{lstlisting}

\subsection{Types of SQL Injection}

\begin{table}[H]
\centering
\caption{SQL injection types}
\begin{tabularx}{\textwidth}{>{\bfseries}lX}
\toprule
Type & Description \\
\midrule
Classic (In-band) & Error or query result returned directly in response \\
Blind Boolean & No visible output — true/false conditions reveal data \\
Blind Time-Based & Response delays reveal data (\texttt{SLEEP(5)}) \\
UNION-Based & UNION statement extracts data from other tables \\
Error-Based & Error messages leak database content \\
\bottomrule
\end{tabularx}
\end{table}

\section{Burp Suite Testing Process}

\subsection{Identifying Injection Points}
\begin{enumerate}
  \item Browse application with Burp proxy on
  \item In HTTP History, identify requests with parameters: \texttt{id=}, \texttt{search=}, \texttt{user=}
  \item Send to Repeater
  \item Add a single quote \texttt{'} to the parameter value
  \item If application returns a SQL error $\to$ injection point confirmed
\end{enumerate}

\begin{lstlisting}[style=http, caption={SQL injection test in Repeater}]
# Original request:
GET /search?q=laptop HTTP/1.1

# Modified (inject single quote):
GET /search?q=laptop' HTTP/1.1

# Vulnerable response:
HTTP/1.1 500 Internal Server Error
{"error": "You have an error in your SQL syntax near 'laptop'' at line 1"}
\end{lstlisting}

\section{Module Completion Checklist}

\begin{tcolorbox}[findingstyle, title={Level 04 Burp Checklist}]
\begin{itemize}
  \item[\checkbox] All parameters with user input identified
  \item[\checkbox] Single quote test applied to each parameter
  \item[\checkbox] SQL error responses identified and documented
  \item[\checkbox] UNION SELECT attempted to identify column count
  \item[\checkbox] Blind injection tested with true/false conditions
  \item[\checkbox] Database version extracted from error or UNION
  \item[\checkbox] All findings documented with full request/response evidence
\end{itemize}
\end{tcolorbox}
